\section{Related Work}\label{sec:related}

Timed automata were introduced by Alur and Dill, who established that reachability is
decidable via the region graph and thereby made automatic verification of real-time systems
possible at all. The region construction is exponential in the number of clocks and too
coarse to implement directly, so practical tools use \emph{zones}, represented as Difference
Bound Matrices, together with abstractions that keep the zone graph finite. The dominant such
abstraction is the $\aLU{}$ family of Behrmann et al., which bounds each clock separately by
the largest constants it is compared against; a complementary line of work represents zones by
minimal constraint systems to reduce storage. These techniques determine the \emph{size} of
the state space and are orthogonal to how a checker traverses it.

\textsc{Uppaal} is the reference implementation and remains the tool against which timed
reachability is measured. It is single-threaded in the configuration we compare against, and
its state counts serve throughout this paper as a correctness oracle: agreement on the verdict
is weak evidence, whereas agreement on the number of symbolic states explored is strong.

Parallel and distributed model checking is well developed for untimed systems, where the
visited set is a plain hash set and the engineering question is a scalable concurrent hash
table. \ltsmin{} is the mature representative, providing a language-independent interface
(PINS) behind which many formalisms plug in, together with a multi-core backend built on a
shared lock-free table. Timed reachability does not fit that mould, because membership is
subsumption rather than equality.

The closest system to ours, and our principal baseline, is \textsc{opaal}+\ltsmin{}, which
generates a successor generator from a \textsc{Uppaal} model and drives it through \ltsmin{}'s
multi-core backend. Its visited set is a shared multimap in which each discrete configuration
owns a list of symbolic states protected by a per-state spinlock, explored either by
work-stealing or by strict level-synchronous BFS. That work reports near-linear speed-ups on
several models and, importantly for us, identifies the failure mode: the size of its critical
regions ``depends crucially on the $|\Sigma|/|\mathcal{S}|$ ratio'', and a high ratio ``causes
low scalability''. It also reports a non-blocking variant that improves scalability at the
cost of allowing $|\Sigma|$ to grow again.

\paragraph*{Positioning.}
We adopt the shared visited set of \textsc{opaal}+\ltsmin{} and depart from it in three ways,
each targeting one of the difficulties in \cref{sec:intro}. Where their expansion routes
successors through per-thread buffers into a locked multimap, we \emph{fuse} expansion and
insertion so a successor is stored by its producer, eliminating the routing phase and the
duplicate copy of each search level that it requires. Where they store a zone per symbolic
state, we intern zones and pack the discrete configuration into a recoverable key, which is
what turns a $13.5$\,GB footprint into $1.0$\,GB. And where they diagnose the
$|\Sigma|/|\mathcal{S}|$ ratio as a limit on scalability and mitigate it by weakening the
invariant, we attack the critical section itself: our batched engine performs the same
subsumption work outside any lock, so the ratio ceases to control contention while the visited
set remains exactly the antichain it was. Finally, whereas a parallel checker that explores
more states than \textsc{Uppaal} can lose to it outright, our merge is aggressive enough that
we match \textsc{Uppaal}'s counts exactly on the models where both are directly comparable.
